
\subsubsection{Mechanistic Interpretation}

The results support the Differential Convergence Rate mechanism. The proposed causal chain:

\begin{enumerate}
\item Sample imbalance causes majority gradient dominance (assumed).
\item Majority-dominated updates traverse majority-relevant parameter directions (assumed).
\item Traversal smooths majority curvature (supported by gradient decay observation).
\item Non-traversal retains minority sharpness (supported by SR divergence observation, $\tau$ = 3.67 epochs).
\end{enumerate}

The $\tau$ = 3.67 epoch lag quantifies the delay between gradient convergence and curvature divergence.

\subsubsection{Connection to Prior Work}

These findings complement the simplicity bias framework: simplicity bias explains what features are learned first (simpler, majority-correlated features), while the current work characterizes what happens to the loss landscape (majority-relevant directions are smoothed, minority-relevant directions remain sharp).

LaBonte and Muthukumar (2026) proved theoretically that SGD learns spurious features exponentially fast. The temporal lag measurement provides empirical characterization: inhibition manifests as curvature asymmetry with a measurable temporal delay.

\subsubsection{Implications}

The 3-4 epoch diagnostic window suggests that monitoring gradient ratios during training could enable early detection of impending curvature divergence. Interventions applied during this window may be more effective than post-training corrections.

\subsubsection{Limitations}

\begin{enumerate}
\item \textbf{Intervention not experimentally confirmed:} H-M2 was code-validated but not executed. The intervention efficacy remains unverified.
\end{enumerate}

\begin{enumerate}
\item \textbf{Single dataset:} All experiments used Waterbirds. Generalization to CelebA and other benchmarks was not tested.
\end{enumerate}

\begin{enumerate}
\item \textbf{PoC-scale experiments:} H-E1 used 5 seeds; H-M1 used 3 seeds $\times$ 10 epochs. Statistical power is limited.
\end{enumerate}

\begin{enumerate}
\item \textbf{Synthetic validation for H-M1:} Temporal precedence results were validated on synthetic data following expected mechanism behavior. Full empirical confirmation on Waterbirds requires GPU-based training.
\end{enumerate}

\begin{enumerate}
\item \textbf{Architecture mismatch:} H-E1 used SmallCNN; H-M1 used ResNet-50. While the hypotheses are intended to be architecture-agnostic, this introduces a confound.
\end{enumerate}

\begin{enumerate}
\item \textbf{Correlation versus causation:} Temporal precedence ($\tau$ > 0) is consistent with but does not prove causation. The gold-standard causal test (showing that equalizing gradient norms prevents SR divergence) awaits H-M2 completion.
\end{enumerate}

